\section{The four-input precode and its implementation}\label{sec:small-scans}

This appendix analyzes the two-scan construction with the four-input
precode, an $[8,4,5]$ MDS code, in place of the Newton precode. With
coefficients from $\{1,\ldots,2048\}$ it proves relative distance above
$14.5\%$ at rate $1/2$, with setup failure below $2^{-56.90}$. The
Newton precode $B_8$ of Section~\ref{sec:width-coefficient} supersedes
this family: with $q=4$ it already proves $18.625\%$
(Table~\ref{tab:width-one-small}), and at each common alphabet size
$q\in\{32,64,128,256\}$ its proved distance is more than twice as large.
For $q\in\{64,128,256\}$, separate timing runs on the same machine
measured the Newton encoder about one millisecond slower
(Tables~\ref{tab:width-one-small} and~\ref{tab:small-timing}).

The encoder retains the two-scan architecture of BAA~\cite{baa26}.
Its random coefficients multiply the running scan states. The analysis
counts support patterns and the ranks of their zero equations, without
taking a union over field-valued messages.

\subsection{Construction and finite guarantee}

Fix a prime $p>q\ge 32$, let $k=2^{20}$ and $n=2k$, and work over $\F_p$.
Partition the input into blocks $(a,b,c,d)$ of four symbols. For each block,
define
\[
 f(t)=a+bt+c\binom{t}{2}+d\binom{t}{3},\qquad
 B(a,b,c,d)=(f(0),\ldots,f(6),d).
\]
Finite differences evaluate this map with 15 additions. Any four output
coordinates determine the input: four evaluations determine a cubic;
three evaluations together with $d$ determine its remaining coefficients.
Thus the four-input precode $B$ is an $[8,4,5]$ MDS code. Apply $B$
independently to all blocks, obtaining a vector of length $n$.

At setup, sample two independent uniform permutations $\pi_1,\pi_2$ of
$\{0,\ldots,n-1\}$. Independently sample each $a_{j,i}$ uniformly from
$\{1,\ldots,q\}$, for $j\in\{1,2\}$ and $1\le i<n$. Round $j$ maps its
input $v$ to $y$ by
\begin{equation}\label{eq:small-scan}
 u_i=v_{\pi_j(i)},\qquad y_0=u_0,\qquad
 y_i=u_i+a_{j,i}y_{i-1}\quad(1\le i<n).
\end{equation}
Output the complete second-round vector. Each round is invertible, so
the code has dimension $k$ for every setup. All vectors after $B$ have
length $n$; there is no systematic copy or final compression.

\begin{theorem}[Finite small-coefficient guarantee]\label{thm:small-scans}
For each fixed prime $p>q$, the construction above satisfies the
following bounds over its setup randomness:
\[
\begin{array}{c@{\qquad}r@{\qquad}c}
 q & d_{\min}\text{ at least}
   & \Pr[\text{failure}]\text{ less than}\\ \hline
 32 & 209716 & 2^{-24.85}\\
 40 & 230897 & 2^{-44.11}\\
 64 & 251659 & 2^{-45.65}\\
 128 & 277873 & 2^{-47.87}\\
 256 & 288359 & 2^{-50.24}\\
 512 & 298845 & 2^{-52.46}\\
 1024 & 301990 & 2^{-54.72}\\
 \mathbf{2048} & \mathbf{304088} & \mathbf{2^{-56.90}}
\end{array}
\]
Failure means that the minimum distance is below the listed threshold.
On a successful setup the distance bound holds for every nonzero message,
including messages selected after the setup is revealed. Failure
exponents are truncated toward zero, so the stated bounds are
conservative.
\end{theorem}

The $q=32$ row uses the sparse bound with $\tau=1$, which ignores
coefficient randomness for source weights at most 128
(Appendix~\ref{app:small-certificate}); this alphabet was not implemented
or timed. Every other row retains coefficient randomness in the
sparse-support calculation.

For $q=40$ the relative distance exceeds $11.01\%$, above the binary
rate-$1/2$ Gilbert--Varshamov (GV) benchmark $H_2^{-1}(1/2)\approx0.110028$.
This comparison is not a binary-code guarantee, and it is far below the
prime-field GV distance. The theorem fixes the prime before setup; it
does not claim that one sampled circuit works simultaneously over primes.
The larger alphabets retain the same precode, two scans, and vector
lengths. Sizes 128, 512, and 2048 give relative distances above
$13.25\%$, $14.25\%$, and $14.5\%$, respectively. All rows from
$q=40$ through $q=2048$ use the same measured scan kernel and memory
footprint. A separate certificate for $q=2048$ gives minimum distance
at least 230897 with failure below $2^{-58.45}$. These are achieved
finite bounds; the table does not claim optimal distance or failure
probability for an alphabet.

Multiplication by a positive integer $m$ has a binary double-and-add
circuit of length $\lfloor\log_2 m\rfloor+\operatorname{popcount}(m)-1$.
For $1\le m\le40$, these lengths sum to 205 and have maximum eight. Including $B$ and the
input additions in~\eqref{eq:small-scan}, the $q=40$ encoder therefore
uses at most $19.875n-18$ additions, with expectation $14.125n-12.25$.
The expectation is over setup, not over messages. The coefficient alphabet
may instead be any $q$ distinct nonzero field elements for the distance
proof; the addition-cost claim uses the positive integer alphabet.
For the measured $q=2048$ configuration, the expected addition
count is $(15815/512)n-14855/512$, about $30.889n$, and the worst-case
count is $(351/8)n-42$. These circuit counts grow with the alphabet.
The native small-integer multiply/add implementation below uses the same
arithmetic instruction sequence for all measured alphabets.

\subsection{Equation ranks instead of a message union}

Fix a source support $S$ and support patterns for both rounds, and condition
on permutations that map these supports as prescribed. The shortened
source space $V_S$ of inputs satisfying the source zero coordinates has
dimension $D_S$. Expose coefficients in round order. Let
$W\subseteq V_S^*$ be the span of all zero equations exposed so far.
A branch remains alive only if $\ker W$ contains a message satisfying
every required nonzero condition whose form is already known.

A \emph{return} is a zero output that ends a run of nonzero outputs.
Its input must be nonzero. The new zero equation has the form
$\ell+a r=0$, where $\ell$ is the current input form and $r$ the
preceding state form. Both are known before exposing $a$. On an alive
branch, $\ell\notin W$. At most one coefficient can make
$\ell+a r\in W$: two distinct such coefficients would imply
$r\in W$ and then $\ell\in W$.

For $\tau\ge1$, put $g=1/q+(1-1/q)/\tau$. Give an alive branch the
\emph{rank potential} $\tau^{-\rank W}$ and a dead branch potential
zero. Each return multiplies the conditional expected potential by at
most $g$; other steps cannot increase it. A realizable nonzero path
has $\rank W\le D_S-1$. Consequently, for a path with $c$ returns,
\begin{equation}\label{eq:small-rank}
 \Pr[\text{prescribed support path is realizable}]
 \le \min\{1,\tau^{D_S-1}g^c\}.
\end{equation}
This argument is adaptive across both rounds. It requires no independence
among the zero equations and incurs no factor $p^{D_S}$.

For a local output support of size $s\in\{5,6,7,8\}$, the source
dimension is $s-4$. The dimension-weighted support polynomial is therefore
\begin{equation}\label{eq:small-source}
 Q(z,\tau)=1+56z^5\tau+28z^6\tau^2+8z^7\tau^3+z^8\tau^4.
\end{equation}
Let $A_g(s,h)$ count admissible input/output support pairs of one scan,
with weights $s,h$, assigning a factor $g$ per return. It is the
coefficient of $x^s y^h$ in
\begin{equation}\label{eq:small-transfer}
 [1,0]\begin{bmatrix}1&xy\\gx&(1+x)y\end{bmatrix}^{n}
 \begin{bmatrix}1\\1\end{bmatrix}.
\end{equation}
The initial state is zero and either final state is allowed. This count
retains support patterns that no field values realize.

For a forbidden output weight $e$, write
$H_g(h,e)=\sum_{j=1}^e A_g(h,j)$ and $L=n/8$.
A uniform permutation maps a fixed support of size $s$ to a uniform
$s$-subset, so each prescribed placement of the source support has
probability $1/\binom ns$, and each placement of the intermediate support
has probability $1/\binom nh$. Combining these normalizations
with~\eqref{eq:small-rank} gives
\begin{equation}\label{eq:small-union}
 \Pr[d\le e]\le\frac1\tau\sum_{s=1}^{4e}
 \frac{[z^s]Q(z,\tau)^L}{\binom ns}
 \sum_{h=1}^{2e}\frac{A_g(s,h)H_g(h,e)}{\binom nh}.
\end{equation}
The inverse scan is bidiagonal, giving $s\le2h\le4e$. A separate $\tau$
may be used for each source weight. Appendix~\ref{app:small-certificate}
gives the finite evaluation of this bound, including the sparse range.

\subsection{Implementation cost}

We implemented alphabets $q\in\{40,64,128,256,512,1024,2048\}$ over
$p=2^{128}-159$. The optimized precode
kernel evaluates eight blocks together with AVX-512, postpones reduction
in 43-bit limbs, and uses contiguous loads and stores. Each scan packs
its 21-bit permutation index and eleven-bit coefficient-minus-one into
one 32-bit word. Thus $q=2048$ is the largest alphabet in this layout.
The scan loop gathers 16 entries at a time and prefetches 256 entries
ahead. Two length-$n$ field buffers suffice; one holds the final output.

\begin{table}[t]
\centering
\caption{Complete encoding at $k=2^{20}$, $n=2^{21}$ and
$p=2^{128}-159$. Times are in milliseconds on one core of an
AMD Ryzen 9 7950X.}
\label{tab:small-timing}
\begin{tabular}{lrrrr}
\toprule
$q$ & Precode $B$ & Scan 1 & Scan 2 & Total\\
\midrule
40 & 1.919 & 7.307 & 7.730 & 16.947\\
64 & 1.924 & 7.317 & 7.741 & 16.989\\
128 & 1.918 & 7.277 & 7.661 & 16.864\\
256 & 1.919 & 7.315 & 7.737 & 16.964\\
512 & 1.922 & 7.337 & 7.739 & 17.009\\
1024 & 1.920 & 7.345 & 7.743 & 17.007\\
\textbf{2048} & \textbf{1.921} & \textbf{7.299} & \textbf{7.704} & \textbf{16.909}\\
\bottomrule
\end{tabular}
\end{table}

Table~\ref{tab:small-timing} reports, for each alphabet, the median of
six process medians, each over 31 timed encodes after two warm-up
encodes, with GCC 15.2 at \texttt{-O3 -march=native}; setup is not timed.
Stage medians need not sum to the median total. All seven totals lie
within $0.5\%$ of the $q=40$ total, and process medians range from
16.855 to 17.075 ms for $q=40$ and from 16.866 to 17.056 ms for
$q=2048$. The larger alphabets therefore show no measurable runtime
penalty.

The optimized implementation uses native small-integer multiplication;
it evaluates exactly the same field map as the addition circuit. The
measured code is therefore distinct from a strict addition-only
implementation, although both have the same distance theorem. Retained
input, work buffers, and metadata occupy
96.008 MiB, excluding validation temporaries and stack scratch.
